\scanpage{28}
Alors la suite \unsure{spectral|spectrale} implique

\[
\begin{aligned}
H^i_Y(X,\unsure{\mu_{nX}|\mu_{n_X}}^{\otimes d})&=0\qquad\text{si }i<2d\\
H^{2d}_Y(X,\mu_{nX}^{\otimes d})&\simeq\Gamma(Y,\underline H^{2d}_Y(\deleted{\unsure{\mathbb Z|Z|\mathbb Z/n\mathbb Z|\mathbb Z/n}}\mu_{nX}^{\otimes d}))
\end{aligned}
\]

d'où un \deleted{\unsure{lieu|hom|base|bon} canonique} élément canonique

\[
\gamma^{(n)}_{X/Y}\in H^{2d}_Y(X,\mu_{nX}^{\otimes d})
\]

[qu'on peut maintenant spécialiser dans $H^{2d}_{\deleted{\unsure{Y|Z}}}(X,\mu_{nX}^{\otimes d})$, ou si $X$ est propre sur un $S$, dans

\[
H^{2d}_{\Phi_{Y/X}}(X,\mu_{nX}^{\otimes d}).
\]

\margin{$\Phi\longrightarrow$ famille de \unsure{spts|cpts}}

\margin{N.B \inserted{\unsure{$d=1$|si|div.}} $\beta_{Y/X}\in H^1_Y(X,G_{mX})$ et $\gamma^{(n)}_{Y/X}=\partial_n(\beta_{Y/X})$, en particulier son image dans $H^2_Y(X,G_{mX})$ est nulle}

\uline{Exemple 1}. $X$ lisse sur $k$ \inserted{sép. clos.}, \deleted{$x$ pt} dim relative $d$, \deleted{il va résulter des le th. de} \inserted{$x$ pt de $X$ rat sur $k$, d'où}

\[
\gamma_{x/X}\in H^{2d}_!(X,\mu_{nX}^{\otimes d}).
\]

Il va résulter de la théorie qu'on va rappeler que \deleted{pour \unsure{tt}} si $X$ est géom. connexe, cet élément est \uline{une base} \uline{de} $H^{2d}_!(X,\mu_{nX}^{\otimes d})$ \uline{et ne dépend pas du choix de $x$}\unsure{;|.} \deleted{\unsure{ce|que} \illegible{1}} C'est qu'on avait vérifié directement dans le cas où $d=1$, en \deleted{\unsure{le|en|les|la} \unsure{compa|compar|corps}}\inserted{nous ramenant} au cas où $C$ était complète et où on dispose \deleted{\unsure{déjà}}

% archive sheet number: 27
